\documentclass[aspectratio=169,12pt]{beamer}
\usepackage{amsmath}
\usepackage{amssymb}
\usepackage{array}
\usepackage[most]{tcolorbox}
\usepackage{graphicx}
\usepackage[sfdefault,lf]{FiraSans}
\renewcommand{\familydefault}{\sfdefault}
\setlength{\parindent}{0pt}
\everymath{\displaystyle}

\setbeamertemplate{navigation symbols}{}
\setbeamertemplate{footline}{}
\setbeamertemplate{headline}{}
\setbeamertemplate{blocks}[rounded][shadow=false]

\definecolor{mmpageWhite}{HTML}{FCFBF7}
\definecolor{mmtanSoft}{HTML}{F7F5EF}
\definecolor{mmgreenDeep}{HTML}{244B2C}
\definecolor{mmgreenLeaf}{HTML}{5D8A56}
\definecolor{mmgreenSoft}{HTML}{E4EFD9}
\definecolor{mmtext}{HTML}{163221}
\definecolor{mmwarn}{HTML}{8F2F36}
\definecolor{mmwarnSoft}{HTML}{F8E8EA}

\setbeamercolor{background canvas}{bg=mmpageWhite}
\setbeamercolor{normal text}{fg=mmtext,bg=mmpageWhite}

\newtcolorbox{MMCard}[2][]{enhanced,breakable=false,colback=#2,colframe=black,boxrule=1.5pt,arc=8pt,left=5pt,right=5pt,top=4pt,bottom=4pt,before skip=0pt,after skip=0pt,#1}
\newcommand{\KoruIcon}{\raisebox{-0.2em}{\includegraphics[height=1.05em]{../../../manamaths/SITE/header-logo.png}}}
\newcommand{\NotesTitle}[1]{{\colorbox{mmtanSoft}{\parbox{0.975\linewidth}{\vspace{0.14em}\hspace{0.25em}{\LARGE\bfseries\textcolor{mmgreenDeep}{#1}\hfill\KoruIcon}\vspace{0.14em}}}\par\vspace{0.18em}{\color{mmgreenLeaf}\rule{\linewidth}{2.0pt}}\par\vspace{0.12em}}}
\newcommand{\SectionTitle}[1]{{\large\bfseries\textcolor{mmgreenDeep}{#1}}\par\vspace{0.04em}}
\newcommand{\GreenDivider}{\par\vspace{0.01em}{\color{mmgreenLeaf}\rule{\linewidth}{2.4pt}}\par\vspace{0.03em}}
\newcommand{\KeyIdea}[1]{\begin{MMCard}[title={\textbf{Key idea}},colbacktitle=mmgreenSoft,coltitle=mmgreenDeep]{mmtanSoft}\small #1\end{MMCard}}
\newcommand{\WorkedExample}[2]{\begin{MMCard}[title={\textbf{#1}},colbacktitle=mmgreenSoft,coltitle=mmgreenDeep]{white}\small #2\end{MMCard}}
\newcommand{\CommonMistake}[1]{\begin{MMCard}[title={\textbf{Common mistake}},colbacktitle=mmwarnSoft,coltitle=mmwarn]{white}\small #1\end{MMCard}}
\newcommand{\TryThese}[1]{\begin{MMCard}[title={\textbf{Try these}},colbacktitle=mmgreenSoft,coltitle=mmgreenDeep]{mmtanSoft}\small #1\end{MMCard}}

\usepackage{tikz}
\setbeamertemplate{background}{
 \begin{tikzpicture}[remember picture,overlay]
 \draw[black,line width=1.5pt,rounded corners=10pt]
 ([xshift=0.45cm,yshift=-0.45cm]current page.north west)
 rectangle
 ([xshift=-0.45cm,yshift=0.45cm]current page.south east);
 \end{tikzpicture}
}

\begin{document}

\begin{frame}[t]
\NotesTitle{Notes \& Steps}
\footnotesize
\KeyIdea{Do the inverse (opposite) to both sides. Whatever you do to one side, do to the other.}
\vspace{-0.2em}
\begin{columns}[T,onlytextwidth]
\begin{column}{0.48\textwidth}
\SectionTitle{Inverse operations}
\vspace{-0.2em}
\begin{align*}
+x &\;\leftrightarrow\; -x \\
\times x &\;\leftrightarrow\; \div x
\end{align*}
\vspace{-0.2em}
\SectionTitle{Solve: \(x+5=12\)}
\vspace{-0.4em}
\begin{align*}
x+5 &= 12 \\
x+5-5 &= 12-5 \\
x &= 7
\end{align*}
\vspace{-0.2em}
Check: \(7+5=12\;\checkmark\)
\end{column}
\begin{column}{0.48\textwidth}
\SectionTitle{Solve: \(3x=15\)}
\vspace{-0.4em}
\begin{align*}
3x &= 15 \\
\frac{3x}{3} &= \frac{15}{3} \\
x &= 5
\end{align*}
\vspace{-0.2em}
Check: \(3\times5=15\;\checkmark\)

\vspace{0.5em}
\SectionTitle{Solve: \(\frac{x}{4}=7\)}
\vspace{-0.4em}
\begin{align*}
\frac{x}{4} &= 7 \\
\frac{x}{4}\times4 &= 7\times4 \\
x &= 28
\end{align*}
\vspace{-0.2em}
Check: \(28\div4=7\;\checkmark\)
\end{column}
\end{columns}
\GreenDivider
\CommonMistake{Forgetting to do the same to both sides. \(x+3=7\) means subtract 3 from both sides, not just move the 3.}
\end{frame}

\begin{frame}[t]
\NotesTitle{Notes \& Steps}
\footnotesize
\begin{columns}[T,onlytextwidth]
\begin{column}{0.48\textwidth}
\SectionTitle{Solve: \(x-3=8\)}
\vspace{-0.4em}
\begin{align*}
x-3 &= 8 \\
x-3+3 &= 8+3 \\
x &= 11
\end{align*}
Check: \(11-3=8\;\checkmark\)

\vspace{1em}
\WorkedExample{Remember:}{
\begin{itemize}
  \item \(+\) undone by \(-\) 
  \item \(-\) undone by \(+\) 
  \item \(\times\) undone by \(\div\) 
  \item \(\div\) undone by \(\times\)
\end{itemize}}
\end{column}
\begin{column}{0.48\textwidth}
\TryThese{\begin{enumerate}
  \item \(x+7=15\)
  \item \(x-5=9\)
  \item \(4x=20\)
  \item \(\frac{x}{3}=6\)
\end{enumerate}}
\vspace{0.5em}
\CommonMistake{Mixing up inverse. For \(x-4=10\) you \textit{add} 4 (not subtract 4). Undo what's been done to \(x\).}
\end{column}
\end{columns}
\end{frame}

\end{document}
